%% =====================================================================
%%  c12：原文印刷页 96–106（PDF 第 97–107 页）
%%  覆盖 3.3.2--3.3.3、3.4--3.4.1、4.--4.1.1，公式 (3.70)--(3.94)、(4.1)--(4.31)
%%  插图：原文 Fig. 11、Fig. 12（原 PDF 第 104 页上两者均只有空白占位框）
%% =====================================================================

\begin{enumerate}[start=2]
  \item 我们将看到，由 \eqref{eq:3.65} 得到的可能的规范不变泛函 $\Psi[A_z]$ 构成一个有限维向量空间。它典范地就是 $\Sigma_2$ 上 WZW 模型的（真空）共形块（conformal block）构成的向量空间。这个空间依赖于 $\Sigma_2$ 上复结构的选择，但同样由于能动张量的 Sugawara 形式，模空间上存在一个射影平坦联络（projectively flat connection）。上述论证最早见于 [17]；更严格的版本在 [7] 中给出。
\end{enumerate}
% (3.65) 位于上一片段 c11（印刷页 95），其 \label{eq:3.65} 应由该片段给出；若编译时此处显示为「??」，请在 c11 中补标签。

\subsubsection{带穿孔曲面上的量子化} % 原文: 3.3.2 Quantization On Punctured Surfaces

我们现在把沿时间方向伸展的 Wilson 线也纳入 $M_3=\Sigma_2\times\mathbb{R}$。它们位于 $\Sigma_2$ 上的点 $P_i\in\Sigma_2$ 处，取 $G$ 的不可约表示，记作 $V(\lambda_i)$，其中 $\lambda_i$ 是与 $P_i$ 处那条线相联系的最高权。

{\small\itshape ［原文标记：说明如何用余伴随轨道的量子力学来表示 Wilson 线。这可以从 Chern-Simons 作用量出发、通过一个奇异的规范变换得到——也可以用别的办法。］}
% [原文待办标记，原文 p.96，页面上以一行星号框出]

我们考虑相空间
\begin{equation}
  \mathcal{P}(E;(\lambda_i,P_i))\equiv\mathcal{A}(E\to\Sigma_2)\times\prod_i\operatorname{Map}(P_i,\mathcal{O}(\lambda_i))
  \tag{3.70}
  \label{eq:3.70}
\end{equation}
其中 $E\to\Sigma_2$ 是复 Hermite 向量丛，$\mathcal{A}(E\to\Sigma_2)$ 是 $E$ 上酉联络（unitary connection）构成的空间，$P_i$ 是 $\Sigma_2$ 上的点，$\mathcal{O}(\lambda_i)$ 是余伴随轨道（coadjoint orbit）：
\begin{equation}
  \mathcal{O}(\lambda_i)=\{\phi=g\lambda_i g^{-1}\}\subset\mathfrak{g}^*
  \tag{3.71}
  \label{eq:3.71}
\end{equation}
$\mathcal{P}(E;(\lambda_i,P_i))$ 是辛流形，且规范群 $\mathcal{G}$ 以辛的方式作用于其上。规范群作用的矩映射（moment map）是
\begin{equation}
  \mu=k\left[\frac{1}{2\pi}F(A)-\sum_i\mu^{(i)}(\phi(P_i))\eta(P_i)\right]
  \tag{3.72}
  \label{eq:3.72}
\end{equation}
其中 $\mu^{(i)}$ 是各轨道的矩映射：
\begin{equation}
  \langle T^A,\mu^{(i)}(g)\rangle=\operatorname{Tr}T^A g\lambda_i g^{-1}=\operatorname{Tr}T^A\phi
  \tag{3.73}
  \label{eq:3.73}
\end{equation}
而 $\eta(P_i)$ 是 $P_i$ 的 Poincaré 对偶（PD）的 delta 函数代表。

% [Tutor Note: (3.72) 有两处极易出错。一是归一化：(3.72) 中 $F(A)$ 前的因子 $1/2\pi$ 与 Chern-Simons 作用量的约定 $S=\frac{k}{4\pi}\int\operatorname{tr}(A\,dA+\frac{2}{3}A^3)$ 绑定，若把作用量改写成 $\frac{k}{2\pi}\int$ 的形式，则 (3.72) 整体要乘 2，$\mu=0$ 所给出的 $F(A)$ 与 Wilson 线源的配比也随之改变；核对时应先回到 (3.66) 附近的 $k$ 归一化。二是 $\eta(P_i)$ 是 $P_i$ 的 Poincaré 对偶的 delta 函数代表，它需要先在 $\Sigma_2$ 上选定体积形式（通常取复结构诱导的度规）才能写出，故 $\mu$ 依赖于该选择；但 $\mu^{-1}(0)/\mathcal{G}$ 不依赖。]

\begin{definition}[共形块空间；原文此处手写编号为“Definition 1”]
共形块空间（space of conformal blocks）
\begin{equation}
  \mathcal{H}(\mathfrak{g},k;(P_i,\lambda_i))
  \tag{3.74}
  \label{eq:3.74}
\end{equation}
是由辛约化（symplectic reduction）$\mu^{-1}(0)/\mathcal{G}$ 的量子化所得到的量子 Hilbert 空间。
\end{definition}
% 排版说明（非批注）：原文以「Definition 1:」手写编号，该编号只在原文这一节内部有意义；
% 这里改用 definition 环境，编号由 LaTeX 按 section 自动生成，故排版后显示为「定义 3.x」而非「定义 1」，
% 引用时请以「共形块空间」这一名称而非编号为准。另注意原文此处商写作 $\mu^{-1}(0)/G$（用 $G$ 而非 $\mathcal{G}$），
% 而 (3.91) 又写 $\mu^{-1}(0)/\mathcal{G}$；按上下文（作用在 $\mathcal{P}$ 上的是规范群）应统一取 $\mathcal{G}$，译文按 $\mathcal{G}$ 处理。

\textbf{评注：}注意，若 $(A;\phi_i)\in\mu^{-1}(0)$，则
\begin{equation}
  P\exp\oint_{\gamma(P_i)}A=\phi_i(P_i)
  \tag{3.75}
  \label{eq:3.75}
\end{equation}
其中 $\gamma(P_i)$ 是环绕 $P_i$ 的小回路。特别地，$P_i$ 处和乐（holonomy）的共轭类被数据 $\lambda_i\in\mathfrak{g}^*$ 固定。

\paragraph{Kähler 量子化} % 原文: 无编号小标题 Kähler Quantization

现在在 $\Sigma_2$ 上选一个复结构，并设 $k\lambda_i$ 落在 $\mathfrak{g}$ 的权格中。于是 $\mathcal{P}(E;(\lambda_i,P_i))$ 成为 Kähler 流形。把它表示成
\begin{equation}
  \mathcal{A}(E\to\Sigma_2)\times\operatorname{Map}(P_i,\mathcal{O}_{k\lambda_i})
  \tag{3.76}
  \label{eq:3.76}
\end{equation}
更便于操作。{\small\itshape ［原文标记：当心：你要的是 $k$ 还是 $k+h$ ？？？］}
% [原文待办标记，原文 p.97 页边，符号为 \ensuremath{\clubsuit}]

$\mathcal{A}$ 上的复结构来自把联络分解为 $(1,0)$ 与 $(0,1)$ 两部分，其中 $\mathcal{A}^{(0,1)}$ 表示酉联络的 $(0,1)$ 部分构成的空间。

众所周知，余伴随轨道作为流形同构于齐性空间。而且，它们可以表示成紧致群或复群的商：
\begin{equation}
  \mathcal{O}_{k\lambda_i}\cong G/H\cong G_{\mathbb{C}}/P(k\lambda_i)
  \tag{3.77}
  \label{eq:3.77}
\end{equation}
其中 $P(k\lambda_i)$ 是与权 $k\lambda_i$ 相联系的抛物子群（parabolic subgroup）。

从这个表示可以清楚地看出，复化规范群 $\mathcal{G}_{\mathbb{C}}$ 在作用。设 $\bar A$ 是酉联络的 $(0,1)$ 部分。则 $\mathcal{G}_{\mathbb{C}}$ 的作用为
\begin{equation}
  (\bar A;\phi_i)\to(\bar A^h;h(P_i)\phi_i(P_i)h(P_i)^{-1})
  \tag{3.78}
  \label{eq:3.78}
\end{equation}

% [Tutor Note: 作者自己在 (3.76) 页边提出了「$k$ 还是 $k+h$」的疑问，这不是无谓的谨慎：几何量子化的预量子化条件只要求 $k\lambda_i$ 落在权格上，而 WZW 模型里真正起作用的能级要按 $k+h^\vee$ 平移——(3.90) 中 $SU(2)$ 的 $k+2$ 正是 $k+h^\vee$（$h^\vee(SU(2))=2$）。因此 (3.76)--(3.79) 写的 $k\lambda_i$ 与第 4 章及 (3.90) 的约定并不一致；引用这些式子时务必检查作者指的是哪一个能级。]

这一空间的描述与 WZW 理论建立了联系。标准做法是先量子化余伴随轨道。于是剩下一个取值于各表示空间中的全纯波函数：
\begin{equation}
  \Psi[\bar A]\in\prod_i\operatorname{Map}(P_i,V(\lambda_i))
  \tag{3.79}
  \label{eq:3.79}
\end{equation}
{\small\itshape ［原文标记：接下来几段我们要换约定：波函数是 $A_z$ 的函数。］}
% [原文待办标记，原文 p.97 页边，符号为 \ensuremath{\clubsuit}]

其中 $V(k\lambda_i)$ 是与该权相联系的 $\mathfrak{g}_{\mathbb{C}}$ 表示。此外，为了能下降到商上、并成为线丛的一个截面，我们必须有
\begin{equation}
  \Psi[\bar A^{h^{-1}}]=e^{kS[h;0,\bar A]}\otimes\rho_{\lambda_i}(h(P_i))\Psi[\bar A]
  \tag{3.80}
  \label{eq:3.80}
\end{equation}

从辛的观点看，这正是 Gauss 定律 $\mu=0$ 的陈述。

\eqref{eq:3.80} 的解照例由「写出到规范不变态上的投影算符」给出。与上面真空共形块的情形一样，我们得到在点 $P_i$ 处插入处于表示 $V(\lambda_i)$ 中的场 $g$ 的 WZW 路径积分：
\begin{equation}
  \Psi(\bar A)=\int[dg]\,e^{S_k[g]-\frac{\mathrm{i}k}{4\pi}\int_{\Sigma_2}\operatorname{Tr}[g^{-1}\partial g\,\bar A]}\prod_i\rho_{\lambda_i}(g(P_i))\Psi_0[\bar A^g]
  \tag{3.81}
  \label{eq:3.81}
\end{equation}
其中 $\Psi_0$ 是任意「试验函数」。

\subsubsection{\texorpdfstring{$D_2\times\mathbb{R}$}{D2xR} 上的量子化（$D_2$ 为圆盘）：非阿贝尔边缘态} % 原文: 3.3.3 Quantization On D2 x R, With D2 A Disk: Nonabelian Edge States

我们现在考虑第 2.3.1 节边缘态论证的非阿贝尔类比。边界条件必须使得
\begin{equation}
  \int_{\partial D\times\mathbb{R}_t}\operatorname{tr}(\delta A\wedge A)=0
  \tag{3.82}
  \label{eq:3.82}
\end{equation}

我们把空间与时间的微分分离出来，$d=d_s+\partial_0dx^0$，同样把联络分离为 $A=A_s+A_0dx^0$。与阿贝尔情形一样，$A_0$ 乘在 Gauss 定律前面，所以把 $A_0$ 当作 Lagrange 乘子积分掉，就强制 $F(A_s)=d_sA_s+A_s^2=0$。

在圆盘上我们可以解出这个方程：
\begin{equation}
  A_s=U^{-1}d_sU
  \tag{3.83}
  \label{eq:3.83}
\end{equation}
其中 $U:D\times\mathbb{R}\to G$。注意 $U$ 仍可以依赖时间。现在把它代回 Chern-Simons 作用量的其余部分，这一点极富启发性：
\begin{equation}
  \exp\left[\mathrm{i}\frac{k}{4\pi}\int_{D\times\mathbb{R}}\operatorname{tr}(A_s\partial_0A_s)dx^0\right]
  \tag{3.84}
  \label{eq:3.84}
\end{equation}

我们需要下面这个简单的恒等式：
\begin{equation}
  \operatorname{tr}\left[(U^{-1}d_sU)\partial_0(U^{-1}d_sU)\right]dx^0
  =-d\left[\operatorname{tr}(U^{-1}d_sU)(U^{-1}\partial_0U)dx^0\right]+\frac{1}{3}\operatorname{tr}(U^{-1}dU)^3
  \tag{3.85}
  \label{eq:3.85}
\end{equation}
于是，在被求值于平坦规范场 $A=U^{-1}dU$ 上之后，Chern-Simons 作用量的其余部分是
\begin{equation}
  \exp\left\{\mathrm{i}\left[-\frac{k}{4\pi}\int_{\partial D\times\mathbb{R}_t}\operatorname{tr}(U^{-1}\partial_\phi U\,U^{-1}\partial_0U)\,d\phi\wedge dx^0
  +\frac{k}{12\pi}\int_{D\times\mathbb{R}}\operatorname{tr}(U^{-1}dU)^3\right]\right\}
  \tag{3.86}
  \label{eq:3.86}
\end{equation}

这样我们就重新得到了「手征」WZW 作用量——注意它对时间只含一阶导数，因此照例应被看作相空间中路径的作用量。

% [Tutor Note: 从 (3.84) 到 (3.86) 靠的是 (3.85) 的恒等式，其中关键一步是把 (3.85) 右端第一个全微分项 $d[\cdots]$ 搬到作用量边界上：它在 $D\times\mathbb{R}$ 内部积分后只剩 $\partial D\times\mathbb{R}_t$ 上的那一项，这正是 (3.86) 的第一个积分，所以不能把这一项当“全导数”丢掉。系数也要核对：$\frac{k}{12\pi}$ 与通常 WZW 项 $\frac{k}{12\pi}\int\operatorname{tr}(U^{-1}dU)^3$ 一致，而 $\frac{k}{4\pi}$ 是 Chern-Simons 的归一化，两者的因子 3 来自 (3.85) 中 $\frac13$ 与 CS 三次项 $\frac23A^3$ 的配合；抄写时把 $12\pi$ 误作 $4\pi$ 是最常见的错误。]

{\small\itshape ［原文标记：把这个推广到包含一条位于中心（center）的 Wilson 线，从而得到非平凡的单值（monodromy）；并求出 $LG/T$ 上的辛形式。］}
% [原文待办标记，原文 p.98，页面上以一行星号框出]

\subsection{平坦规范场模空间的量子化与代数几何} % 原文: 3.4 Quantization Of The Moduli Space Of Flat Gauge Fields And Algebraic Geometry

也可以先施加约束：即先下降到有限维的平坦联络模空间——把它看成一个辛空间——再量子化它。这条途径会引出一些相当不平凡的数学。

同样，最有用的是给 $\Sigma_2$ 一个复结构。这会在 $\mathcal{M}_{flat}$ 上诱导一个复结构。其机制如下：我们可以把协变导数 $D$ 分解为 $D=D^{1,0}+D^{0,1}$。因为我们处在一个复维，$(D^{0,1})^2=0$。平坦联络空间的切空间可以等同于 $D^{0,1}$ 的核，于是
\begin{equation}
  T^{1,0}\mathcal{M}_{flat}\cong H^{0,1}(E)\cong\ker D^{0,1}/\mathit{guage}
  \tag{3.87}
  \label{eq:3.87}
\end{equation}

于是预量子线丛（prequantum line bundle）是 $\mathcal{L}=\operatorname{DET}(D^{0,1})$，即行列式线丛（determinant line bundle）。

按 Narasimhan-Seshadri 的一个重要定理，全纯空间 $\mathcal{M}_{flat}$ 可以等同于某个全纯向量丛的模空间。

% 笔误说明（非批注）：(3.87) 分母上的 “guage” 是原文的笔误（应为 gauge），此处照原文转写、未作改动；
% 它指的就是按规范群 $\mathcal{G}$ 取商。

对 $G=SU(N)$，我们考察秩为 $N$、行列式固定的全纯「半稳定」向量丛。{\small\itshape ［原文标记：此处需要解释］}
% [原文待办标记，原文 p.99]

$\mathcal{M}_{holvb}$ 与某个射影空间双有理等价 {\small\itshape ［原文标记：要解释］}，因此 $\operatorname{Pic}(\mathcal{M}_{holvb})\cong\mathbb{Z}$。其生成元可以相当漂亮地描述如下：取 Riemann 面 $\Theta$-除子的一个元素，即取 $L\in\operatorname{Pic}^{g-1}(\Sigma_2)$ 使 $h^0(L)\neq0$。然后定义
\begin{equation}
  \Theta_L\subset\mathcal{M}_{holvb}
  \tag{3.88}
  \label{eq:3.88}
\end{equation}
为使得 $h^0(E\otimes L)\neq0$ 的那些丛 $E$ 构成的集合。可以证明 {\small\itshape ［原文标记：文献待补］}，对 $\Sigma_2$ 的 $\Theta$-除子中的两个选择 $L,L'$，有 $\Theta_L\sim\Theta_{L'}$，于是我们典范地得到 $\mathcal{M}_{holvb}$ 中的一个除子。

事实上，$\operatorname{Pic}(\mathcal{M}_{holvb})$ 的生成元同构于行列式线丛。

基本想法：给 $E$ 一个全纯联络，于是我们要找的是使 $\bar\partial_Es=0$ 有非零解的丛，而这正是 $\operatorname{DET}(\bar\partial_E)$ 的除子。

对代数几何学家而言，共形块空间于是就是
\begin{equation}
  H^0(\mathcal{M}_{hol.v.b.};\mathcal{L}^k)
  \tag{3.89}
  \label{eq:3.89}
\end{equation}
其中 $\mathcal{L}$ 生成 Picard 群。

我们将会看到，从共形场论出发来思考，会相当自然地导出一个关于共形块空间维数的著名公式，即 Verlinde 公式。对 $G=SU(2)$、能级为 $k$、亏格为 $g$ 的情形，它有如下显式形式：
\begin{equation}
  \dim H^0(\mathcal{M}_{hol.v.b.};\mathcal{L}^k)
  =\left(\frac{k+2}{2}\right)^{g-1}\sum_{n=1}^{k+1}\left(\sin\left(\frac{n}{k+2}\pi\right)\right)^{2-2g}
  \tag{3.90}
  \label{eq:3.90}
\end{equation}

% [Tutor Note: (3.90) 值得做一个小核对：$g=1$ 时前因子 $((k+2)/2)^{0}=1$，求和号给出 $k+1$，正是 $SU(2)$ 能级 $k$ 的初级场（primary field）个数，也就是手征代数的不可约表示数——环面配分函数确有 $k+1$ 个特征标，故这一步可以放心。另注意式中 $k+2$ 应读作 $k+h^\vee$（$h^\vee(SU(2))=2$），这正是 (3.76) 页边那条「$k$ 还是 $k+h$」疑问的答案：Verlinde 公式里出现的是平移后的能级。$g\ge2$ 时指数 $2-2g<0$，但 $n\pi/(k+2)\in(0,\pi)$ 使 $\sin>0$，故不涉及符号问题。]

{\small\itshape ［原文标记：关于 $k\to\infty$ 极限的评注。半经典极限，因而要取相空间的辛体积；这可由二维 Yang-Mills 得到：标度 $g=e^{\phi/k}$，得到 $\int\operatorname{Tr}\phi F$。］}
% [原文待办标记，原文 p.99，页面上以星号框出]

{\small\itshape ［原文标记：曲线模空间上的射影平坦联络：在此语境下它被称为「Hitchin 联络」——引用 Hitchin 的文章。要解释吗？］}
% [原文待办标记，原文 p.99]

\subsubsection{含顶点算符插入的共形块的推广} % 原文: 3.4.1 Generalization To Conformal Blocks With Vertex Operator Insertions

等变性条件 \eqref{eq:3.80} 表明，波函数 $\Psi$ 定义了 $[\mathcal{P}(E;(\lambda_i,P_i))]^{stable}/\mathcal{G}_{\mathbb{C}}$ 上某个全纯线丛的截面。上标「stable」指的是这样一个事实：为了得到好的商，我们必须对 $\mathcal{G}_{\mathbb{C}}$ 所作用的规范场空间加以限制——这在几何不变量理论（geometric invariant theory）中是熟悉的事。

与 Narasimhan-Seshadri 定理类似的一个陈述是 Mehta-Seshadri 定理，它把商
\begin{equation}
  \mu^{-1}(0)/\mathcal{G}\cong[\mathcal{P}(E;(\lambda_i,P_i))]^{stable}/\mathcal{G}_{\mathbb{C}}
  \tag{3.91}
  \label{eq:3.91}
\end{equation}
等同于「带抛物结构的稳定丛」的模空间。我们简要说明这指的是什么。

按 Borel-Weil-Bott 定理，我们可以把最高权表示 $\lambda$ 与 $\mathcal{G}_{\mathbb{C}}$ 中的旗（flag）相联系。相应地，我们可以考虑全纯丛 $\mathcal{E}\to X_2$ 的模空间，其中向量丛内部指定了旗：
\begin{equation}
  \mathcal{E}_0\subset\mathcal{E}_1\subset\cdots\subset\mathcal{E}_n
  \tag{3.92}
  \label{eq:3.92}
\end{equation}
旗位于各点 $P_i$ 处。带有这种结构的全纯向量丛称为准抛物（quasi-parabolic）的。我们现在给数据 \eqref{eq:3.92} 配上一列实数，称为「权」$\alpha_i$，其重数由旗决定。（在实的情形，这些数将对应于 \eqref{eq:3.75} 中和乐的共轭类。）\eqref{eq:3.92} 的数据连同这些权一起定义了一个抛物向量丛。

为了定义一个良好的模问题，需要借助正规化次数（normalized degree）来给出抛物向量丛稳定性的定义：
\begin{equation}
  \mathit{par}\text{-}\mathit{deg}(\mathcal{E})=\mu(\mathcal{E})=\frac{c_1(\mathcal{E})}{\mathit{rank}(\mathcal{E})}+\sum\mathit{WEIGHTS}
  \tag{3.93}
  \label{eq:3.93}
\end{equation}

粗略地说，稳定性的定义禁止存在 $\mu$ 更大的全纯子丛。这个条件排除全纯自同构的存在，并保证 \eqref{eq:3.91} 中的商是好的。

我们现在用 \eqref{eq:3.91} 右端的全纯量子化来描述共形块空间。于是，共形块空间就是某个线丛的全纯截面空间：
\begin{equation}
  \mathcal{H}(\mathfrak{g},k;(P_i,\lambda_i))
  =H^0\left\{[\mathcal{P}(E;(\lambda_i,P_i))]^{stable}/\mathcal{G}_{\mathbb{C}};\mathcal{L}^{\otimes k}\right\}
  \tag{3.94}
  \label{eq:3.94}
\end{equation}

\section{WZW 模型与有理共形场论} % 原文: 4. WZW Model And RCFT

\begin{keytakeaways}
  \item 物理图像：二维 WZW 模型 = 主手征模型（度规耦合的 sigma 模型）加一个 Wess-Zumino 项。后者无法写成 $2n$ 维时空上的积分，必须把时空的像「填满」一张 $(2n+1)$ 维链，但它的变分只落在 $2n$ 维边界上，因此仍给出局域的运动方程。
  \item 数学工具：$\Theta=g^{-1}dg$ 是 $G$ 上的 Maurer-Cartan 形式，$\operatorname{Tr}\Theta^{2n+1}$ 闭且（当 $\operatorname{rank}G$ 足够大时）代表非零上同调类；适当归一化的 $x_{2n+1}=[c_n\operatorname{Tr}\Theta^{2n+1}]$ 生成 $H^{2n+1}_{DR}(G)$ 的整系数格，$SU(N)$ 情形可取 $x_3=[\frac{1}{24\pi^2}\operatorname{Tr}_N(g^{-1}dg)^3]$。
  \item 核心结论：换用另一张不与之小形变的包围链，$WZ(g)$ 只会改变 $2\pi kN$（$N\in\mathbb{Z}$），故 $WZ(g)\bmod 2\pi k\mathbb{Z}$ 良定义；再要求 $e^{\mathrm{i}WZ/\hbar}$ 在 $\operatorname{Map}[S_{2n},G]$ 上单值，就得到 $k=\kappa\hbar$——耦合常数的拓扑量子化。
  \item 计算要点：运动方程为 $-\frac{f^2}{2}d(*g^{-1}dg)+2\pi kc_n(2n+1)(g^{-1}dg)^{2n}=0$；WZ 项满足 Polyakov-Wiegman 型的乘法律 \eqref{eq:4.28}，这是把它当作二维共形场论作用量来使用的代数基础。
\end{keytakeaways}

WZW 模型是二维有理共形场论（rational conformal field theory, RCFT）的一个非常重要的例子。它有着庞大而丰富的文献，而且仍有许多有趣的问题有待更好地理解。例如，当 $G$ 为非紧半单 Lie 群时，$LG$ 的理论还远远没有被完全理解。二维量子 WZW 模型的应用包括：

\begin{enumerate}
  \item 非平凡的可精确求解的量子（共形）场论。

  \item 按 Moore + Seiberg 的一个猜想，它们是所有有理共形场论的基础。
  {\small\itshape ［原文标记：关于这个猜想及其当前状况的评论放到别处去］}
  % [原文待办标记，原文 p.100 页边，符号为 \ensuremath{\clubsuit}]

  \item 在弦论中的大量用途：异弦（heterotic string）的规范对称性、可精确求解的弦背景、孤子五膜（solitonic five-brane）、显式的非平凡 D-膜态、AdS/CFT 的可精确求解的实例，等等。

  \item 在凝聚态物理中的大量用途：二维临界现象、Kondo 问题、量子 Hall 效应，以及在量子计算中的应用。

  \item 数学：正如本讲义所解释的，它们给出了模张量范畴（modular tensor category）的一类显式而重要的例子，从而与纽结不变量、三维 Chern-Simons 理论、曲面上平坦联络的模空间建立了联系。此外，它们还在 $K$-理论的一些非常引人注目的结果（Freed-Hopkins-Teleman）中起作用。它们在几何 Langlands 纲领中也扮演重要角色，其中能级 $k=-h$ 的代数具有特殊地位。
\end{enumerate}

% [Tutor Note: 首句的 $LG$ 在本文其余各处均指 $S^1$ 上映射构成的环群（loop group）$\operatorname{Map}(S^1,G)$ 的中心扩张（见第 4 章后文 “$LG$ is a central extension of the loop group $LG=\operatorname{Map}(S^1,G)$” 以及 $(LG)_L\times(LG)_R$ 的记号）。因此「$LG$ 的理论」宜读作「以环群为手征对称代数的 WZW 理论」，而不是「环群上的某个独立理论」；原文此处措辞简省，翻译时按字面保留 $LG$。]

\subsection{主手征模型与 WZ 项} % 原文: 4.1 The Principal Chiral Model And The WZ Term

\subsubsection{Wess-Zumino 项} % 原文: 4.1.1 The Wess-Zumino Term

考虑映射 $g:S_d\to G$ 的 sigma 模型，其中 $G$ 是 Lie 群，$S_d$ 是 $d$ 维的（伪）Riemann 时空。该理论的标准 sigma 模型作用量是
\begin{equation}
  \frac{f^2}{4}\int_{S_d}\operatorname{Tr}(g^{-1}dg)\wedge*(g^{-1}dg)
  \tag{4.1}
  \label{eq:4.1}
\end{equation}
其中 $f$ 是耦合常数，$*$ 是 Hodge 星算子。

对规范场论中反常（anomaly）的考察，促使 Wess 与 Zumino 在四维 sigma 模型的情形下往 sigma 模型作用量里引入了一个非常有趣的项 [63]。它恰当的概念表述与物理后果，随后由 Witten 在一系列文章 [64, 65] 中漂亮地澄清了。我们这里把它写在任意偶数时空维数 $d=2n$ 的情形。

设 $\Theta=g^{-1}dg$ 是 $G$ 上的 Maurer-Cartan 形式。则 $\operatorname{Tr}\Theta^{2n+1}$ 是闭形式。若 $G$ 的秩适当地大于 $n$（我们的主要应用是 $n=1$，此时这总是成立），则它代表一个非零上同调类，并且对适当的归一化常数 $c_n$，
\begin{equation}
  x_{2n+1}=[c_n\operatorname{Tr}\Theta^{2n+1}]
  \tag{4.2}
  \label{eq:4.2}
\end{equation}
是一个 de Rham 上同调类，它生成 $H^{2n+1}_{DR}(G)$ 中的整系数格。

设 $g:S_{2n}\to G$ 是 sigma 模型场，并设时空是闭的，即 $\partial S_{2n}=0$。物理上，即便对 $\mathbb{R}^{2n}$ 上的场，这也是相关的，只要我们要求场在空间与时间无穷远处趋于 $1$。在那种情形下，我们可以认为场定义在 $S^{2n}$ 上。\footnote{推广到时空带边界的情形非常有意思。那时 $\exp[\mathrm{i}WZ]$ 应被看作某个线丛的截面。}

定义 Wess-Zumino 项有几种略有不同的途径。一种做法是注意到像
\begin{equation}
  g(S_{2n})\subset G
  \tag{4.3}
  \label{eq:4.3}
\end{equation}
是 $G$ 内部的一个 $(2n)$-闭链，它随 $G$ 连续变化。现在，若 $H_{2n}(G;\mathbb{Z})=0$（情形常常如此\footnote{例如对紧致、单、单连通的 Lie 群总有 $H_2(G;\mathbb{Z})=0$。除 $\pi_4(USp(2n))=\mathbb{Z}_2$ 之外，对所有紧致、单、单连通的群都有 $\pi_4(G)=0$。{\small\itshape ［原文标记：列出 $H_4(G;\mathbb{Z})$］}}），我们就可以用一张有向链 $\mathcal{B}_{2n+1}(g)$ 把时空的像填满：
\begin{equation}
  \partial\mathcal{B}_{2n+1}(g)=g(S_{2n}).
  \tag{4.4}
  \label{eq:4.4}
\end{equation}

这张链也随 $g$ 连续变化。但请注意，$\mathcal{B}_{2n+1}(g)$ 可以有不同的选择。

\textbf{例}\quad $S=S^2$，$G=SU(2)\cong S^3$，映射 $g$ 把 $S$ 送到赤道。那么我们可用上半球 $D_+$。

我们把 Wess-Zumino 项定义为
\begin{equation}
  WZ(g):=2\pi k\int_{\mathcal{B}_{2n+1}(g)}\omega_{2n+1}
  \tag{4.5}
  \label{eq:4.5}
\end{equation}
其中 $k$ 是一个实的耦合常数，具有 $\hbar$ 的量纲。WZW（Wess-Zumino-Witten）理论就是具有 Minkowski 空间作用量
\begin{equation}
  \frac{f^2}{4}\int_{S_d}\operatorname{Tr}(g^{-1}dg)\wedge*(g^{-1}dg)+WZ(g)
  \tag{4.6}
  \label{eq:4.6}
\end{equation}
的非线性 sigma 模型。

乍看之下，定义 \eqref{eq:4.5} 显得荒谬。有两个直接的麻烦：
\begin{itemize}
  \item 它看起来是 $2n+1$ 维中场构型的作用量。
  \item 它看起来依赖于包围链 $\mathcal{B}_{2n+1}$ 的选择，而约束 \eqref{eq:4.4} 给 $\mathcal{B}_{2n+1}$ 留下了无穷多个选择。
\end{itemize}

让我们先处理第 (1) 点。虽然 WZ 项的定义用到了 $2n+1$ 维的场构型，但作用量的\emph{变分}只依赖于 $2n$ 维边界 $\partial\mathcal{B}_{2n+1}$ 上的场，因此作用量实际上是局域的！关于如何变分作用量并导出运动方程的细节，见下面的练习。

于是，我们得到 WZ 项的变分：
\begin{equation}
  \delta WZ(g)=2\pi kc_n(2n+1)\int_{S_{2n}}\operatorname{Tr}(g^{-1}\delta g)(g^{-1}dg)^{2n}
  \tag{4.7}
  \label{eq:4.7}
\end{equation}

值得注意的是，尽管 WZ 作用量的定义涉及向高一维的延拓，它仍是个局域作用量，意思是：它在场 $g(x)$ 的局域改变下的变分是时空上的局域密度！它的值可能依赖微妙的拓扑问题，但变分是局域的。

因此，WZW 理论的运动方程是局域的偏微分方程：
\begin{equation}
  -\frac{f^2}{2}d(*g^{-1}dg)+2\pi kc_n(2n+1)(g^{-1}dg)^{2n}=0
  \tag{4.8}
  \label{eq:4.8}
\end{equation}

% [Tutor Note: (4.7) 是本片段最容易被“跳过”的一步：作者只说「于是我们得到」，实际用到了 (i) $\delta\Theta=\delta(g^{-1}dg)=d(g^{-1}\delta g)+[g^{-1}dg,g^{-1}\delta g]$（即 (4.29)），(ii) 对 $\operatorname{Tr}\Theta^{2n+1}$ 变分时对易子项的迹为零（因 $\operatorname{Tr}$ 的循环性与不变性），(iii) 把 $2n$ 个因子的 $g^{-1}dg$ 留下、把 $d(g^{-1}\delta g)$ 分部积分到 $S_{2n}$ 上。这三步的完整推演在 p.106 的练习 “Variation Of The WZ Term”（(4.29)--(4.31)）里，系数 $(2n+1)$ 就来自 $\delta$ 落在 $2n+1$ 个因子中的任意一个。]

现在来处理第二点——对包围链 $\mathcal{B}_{2n+1}(g)$ 选择的依赖。对固定的 $g$，我们当然可以把链光滑地形变而得到第二条链，如图 11 所示。差 $\mathcal{B}'-\mathcal{B}$ 是 $G$ 中一个小的闭 $2n+1$ 闭链，而且它同调于零，故 $\mathcal{B}'-\mathcal{B}=\partial\mathcal{Z}$，其中 $\mathcal{Z}$ 是一张 $(2n+2)$-链。但由 Stokes 定理：

% 原 PDF 第 104 页上，Fig. 11 与 Fig. 12 的图形区域为纯白（无任何线条），
% 即作者原稿此处只有空白占位框；figures/fig11.png、figures/fig12.png 的裁切结果
% 因而也是空白（fig12 的裁切还误含了一截图注文字）。此处按 main.tex 中既定的
% \MissingFigureNote 办法处理：保留图题与编号，不插入空白图。
\begin{figure}[htbp]
  \centering
  \MissingFigureNote
  \caption{$G$ 中两条略有差别的 $(2n+1)$-链 $\mathcal{B}$ 与 $\mathcal{B}'$，它们围出同一个 $2n$-闭链 $g(S_{2n})$。（原文 Fig.~11，第 103 页）}
  \label{fig:11}
\end{figure}

\begin{figure}[htbp]
  \centering
  \MissingFigureNote
  \caption{$G$ 中两条不同的 $(2n+1)$-链，它们围出同一个 $2n$-闭链 $g(S_{2n})$。（原文 Fig.~12，第 103 页）}
  \label{fig:12}
\end{figure}

\begin{equation}
  \int_{\mathcal{B}'}\omega_{2n+1}-\int_{\mathcal{B}}\omega_{2n+1}=\int_{\mathcal{Z}}d\omega_{2n+1}=0
  \tag{4.9}
  \label{eq:4.9}
\end{equation}

于是 $WZ(g)$ 在 $\mathcal{B}$ 的小形变下不变。

然而，也可能出现 $\mathcal{B}'$ 与 $\mathcal{B}$ 并非彼此的小形变的情形，如图 12。一般地，若 $\mathcal{B},\mathcal{B}'$ 是两张有向链，满足
\begin{equation}
  \partial\mathcal{B}=\Sigma_{2n}
  \tag{4.10}
  \label{eq:4.10}
\end{equation}
与
\begin{equation}
  \partial\mathcal{B}'=\Sigma_{2n}
  \tag{4.11}
  \label{eq:4.11}
\end{equation}
则
\begin{equation}
  \mathcal{B}\cup-\mathcal{B}'=\Xi_{2n+1}
  \tag{4.12}
  \label{eq:4.12}
\end{equation}
是一个闭的有向 $(2n+1)$-闭链。因此，
\begin{equation}
  \int_{\mathcal{B}(g)}\omega_{2n+1}=\int_{\mathcal{B}'(g)}\omega_{2n+1}+\int_{\Xi_{2n+1}}\omega_{2n+1}
  \tag{4.13}
  \label{eq:4.13}
\end{equation}
从而，若周期 $\int_{\Xi_{2n+1}}\omega_{2n+1}$ 非零，则表达式 $WZ(g)$ 作为实数并未良定义！

这或许令人不安，但闭链 $\Xi_{2n+1}$ 定义了一个整同调类，因此 $\omega_{2n+1}$ 的周期是量子化的。于是 $WZ(g)$ 定义中的含糊性正是形如 $2\pi kN$ 的加性量子化平移，其中 $N$ 是整数。换言之，
\begin{equation}
  WZ(g)\bmod 2\pi k\mathbb{Z}
  \tag{4.14}
  \label{eq:4.14}
\end{equation}

是良定义的。量子化的含糊性不会随 $g$ 的小改变而改变。因此 $WZ(g)$ 仍是局域作用量，运动方程也仍是局域的。

注意此处的状况与我们在讨论辛流形的一般量子化时的作用量、而辛形式具有非平凡周期时的状况非常相似。

量子力学中的情形要更微妙一些，因为在量子力学中人们直接用作用量、而不只是运动方程。不过在量子力学中，作用量只通过 $\exp[\frac{\mathrm{i}}{\hbar}S]$ 进入，因此真正必须良定义的只是表达式
\begin{equation}
  \exp\left[\frac{\mathrm{i}}{\hbar}WZ(g)\right]
  \tag{4.15}
  \label{eq:4.15}
\end{equation}

\eqref{eq:4.15} 中的含糊性是什么？我们看到它正是
\begin{equation}
  \exp\left[2\pi\mathrm{i}\frac{k}{\hbar}\int_{\Xi_{2n+1}}\omega_{2n+1}\right]
  \tag{4.16}
  \label{eq:4.16}
\end{equation}

因此，如果 $k=\kappa\hbar$ 且 $\kappa$ 为整数，那么
\begin{equation}
  \exp\left[\frac{\mathrm{i}}{\hbar}WZ(g)\right]
  \tag{4.17}
  \label{eq:4.17}
\end{equation}
在路径积分中就是场空间 $\operatorname{Map}[S_{2n},G]$ 上良定义的 $U(1)$-值函数。假定场空间上有良定义的测度，把这个表达式放进测度里也无妨。

于是，为使量子力学路径积分中的测度在数学上良定义，耦合常数 $k$ 必须量子化。这是耦合常数的拓扑量子化最漂亮的例子之一。

我们通常取 $\hbar=1$。于是大的 $k$ 对应半经典极限。

\textbf{Remarks（注记）：}

\begin{enumerate}
  \item \emph{整系数归一化}。以下是一些相关事实。可以证明\footnote{参见 Mimura and Toda, \emph{Topology of Lie Groups}, Translations of Math. Monographs \textbf{91}; R. Bott, Bull. Soc. Math. France \textbf{84}(1956) 251.}：对所有紧致、单、连通且单连通的群 $G$，
  \begin{equation}
    x_3=\left[\frac{1}{48\pi^2h^\vee}\operatorname{Tr}_{adj}(g^{-1}dg)^3\right]
    \tag{4.18}
    \label{eq:4.18}
  \end{equation}
  生成 $H^3_{DR}(G)$ 中的整上同调格，其中 $h$ 是对偶 Coxeter 数（dual Coxeter number）。特别地，对 $SU(2)$ 我们可以取
  \begin{equation}
    x_3=\left[\frac{1}{24\pi^2}\operatorname{Tr}_2(g^{-1}dg)^3\right]
    \tag{4.19}
    \label{eq:4.19}
  \end{equation}
  由此可知，对 $SU(N)$ 我们可以取
  \begin{equation}
    x_3=\left[\frac{1}{24\pi^2}\operatorname{Tr}_N(g^{-1}dg)^3\right]
    \tag{4.20}
    \label{eq:4.20}
  \end{equation}
  来生成 $H^3_{DR}(G)$ 中的整上同调格。

  \item 这里是定义 Wess-Zumino 项的另一种方式。对场空间 $\operatorname{Map}(S_{2n},G)=\coprod_{\alpha}C_{\alpha}$ 的每个连通分支 $C_{\alpha}$，我们选一个「基点」场构型 $g_0^{(\alpha)}:S_{2n}\to G$。若 $S_{2n}$ 可缩，则只有一个分支，我们可以取 $g_0$ 为常映射（比如像为 $1\in G$）。一般地，对场构型 $g\in C_{\alpha}$，我们取一条光滑同伦 $g(x,s)$，$0\le s\le1$，使其在 $s=0$ 时为 $g_0^{(\alpha)}(x)$、在 $s=1$ 时为 $g(x)$。现在把这一插值视为 $2n+1$ 维中的场，即柱面 $\hat g:I\times S_{2n}\to G$ 的映射。于是我们可以定义
  \begin{equation}
    WZ(g;g_0):=2\pi k\int_{I\times S_{2n}}\hat g^*(\omega_{2n+1})
    \tag{4.21}
    \label{eq:4.21}
  \end{equation}

  \item $WZ(g;g_0)$ 的值依赖于 $g_0$ 的选择以及插值，但只是「局域地」依赖，其意义如下：设有连通分支 $C_{\alpha}$ 中的一族连续映射 $\tilde g^\tau:S_{2n}\to G$。那么我们可以找到一族连续延拓 $\tilde g^\tau_s:B_{2n+1}\to G$，使得对所有 $\tau$ 都有 $g^\tau_0(x)=g_0(x)$。于是，令 $B=I\times S_{2n}$，我们有
  \begin{equation}
    \frac{\partial}{\partial\tau}\operatorname{Tr}(\tilde g^{-1}d_B\tilde g)^{2n+1}
    =d_B\left[(2n+1)\operatorname{Tr}\tilde g^{-1}\frac{\partial\tilde g}{\partial\tau}(\tilde g^{-1}d_B\tilde g)^{2n}\right]
    \tag{4.22}
    \label{eq:4.22}
  \end{equation}
  % 原文此式外加了方框（作者的强调记号），此处未加框，仅以本条注释说明。

  证明：我们知道，拉回到 $I\times B_{2n+1}$ 上的 Maurer-Cartan 形式是闭的，故
  \begin{equation}
    (d_B+\delta)\operatorname{Tr}(\tilde g^{-1}(d_B+\delta)\tilde g)^{2n+1}=0
    \tag{4.23}
    \label{eq:4.23}
  \end{equation}
  其中 $\delta=d\tau\frac{\partial}{\partial\tau}$。现在，乘积空间上的形式可以按型 $(a,b)$ 分解，其中 $a$-形式沿 $I$、$b$-形式沿 $B_{2n+1}$。\eqref{eq:4.23} 中 $(1,2n+1)$ 型的成分是
  \begin{equation}
    \delta\operatorname{Tr}(\tilde g^{-1}d_B\tilde g)^{2n+1}
    +(2n+1)d_B\left[\operatorname{Tr}(\tilde g^{-1}\delta\tilde g)(\tilde g^{-1}d_B\tilde g)^{2n}\right]=0
    \tag{4.24}
    \label{eq:4.24}
  \end{equation}
  抽出 $d\tau$ 即得我们的恒等式。现在很容易证明，按 \eqref{eq:4.21} 定义的 WZ 项的变分只依赖于 $s=1$ 处 $g^\tau_s$ 的变分。
\end{enumerate}

% [Tutor Note: 「It is now an easy matter to show that ...」一句省掉的是对 (4.22) 沿 $\tau$ 积分并处理边界项的全部工作：因为对一切 $\tau$ 有 $g^\tau_0=g_0$（即 $s=0$ 端固定），而由 (4.22) 右边是全微分 $d_B(\cdots)$，在 $B=I\times S_{2n}$ 上积分后只剩 $s=1$ 端面 $S_{2n}$ 上的贡献；这正是「只依赖 $s=1$ 处的变分」的含义，也说明 (4.21) 与包围链定义 (4.5) 相差的只是一个与 $g_0$ 有关的常数。若把这一句当作「显然」跳过，(4.22) 里那个 $(2n+1)$ 因子的来历以及为何不是 $2n+2$ 就会变得难以核对。]

\medskip\noindent\textbf{练习}
\begin{enumerate}[label=\textup{\alph*.)}]
  \item 用群的 Euler 角、取基本表示，计算 $SU(2)$ 的 $\operatorname{Tr}(g^{-1}dg)^3$：
  \begin{equation}
    \operatorname{Tr}_2(g^{-1}dg)^3=-\frac{3}{2}d\psi\wedge\sin\theta\,d\theta\wedge d\phi
    \tag{4.25}
    \label{eq:4.25}
  \end{equation}
  \item 把这个微分形式写成局部恰当的形式。
  \item 证明
  \begin{equation}
    \int_{SU(2)}\frac{1}{24\pi^2}\operatorname{Tr}_2(g^{-1}dg)^3=-1
    \tag{4.26}
    \label{eq:4.26}
  \end{equation}
  从而断定该形式不是整体恰当的。与上面的一般归一化相比较。
  \item 再证明方程 \eqref{eq:4.20} 中的 $x_3$ 对一切 $SU(N)$ 都定义了非平凡的上同调类。
\end{enumerate}

% [Tutor Note: (4.26) 右端是 $-1$ 而不是 $+1$，这一点必须与 (4.19)、(4.18) 一起读：$[\,\frac{1}{24\pi^2}\operatorname{Tr}_2(g^{-1}dg)^3\,]$ 在 $SU(2)\cong S^3$ 上的积分等于 $-1$，即它生成整上同调格但取负的定向；所以 (4.26) 与 (4.19) 并不矛盾，只差 $SU(2)$ 作为 $S^3$ 的定向约定（等价地，差一个 $g\to g^{-1}$ 的替换）。另外 (4.18) 的正文说「$h$ 是对偶 Coxeter 数」而式中写作 $h^\vee$，两者是同一个量，转写时保留原样。]

\medskip\noindent\textbf{练习}\quad\emph{Polyakov-Wiegman 公式}
% 原文拼作 Polyakov-Wiegman；通行拼法为 Polyakov-Wiegmann，此处照原文保留。

考虑二维时空中的 WZ 项。

\begin{enumerate}[label=\textup{\alph*.)}]
  \item 证明
  \begin{equation}
    \operatorname{Tr}((g_1g_2)^{-1}d(g_1g_2))^3
    =\operatorname{Tr}(g_1^{-1}dg_1)^3+\operatorname{Tr}(g_2^{-1}dg_2)^3
    +3d\left[\operatorname{Tr}(dg_2g_2^{-1})(g_1^{-1}dg_1)\right]
    \tag{4.27}
    \label{eq:4.27}
  \end{equation}
  \item 由此断言 WZ 项满足
  \begin{equation}
    WZ(g_1g_2)=WZ(g_1)+WZ(g_2)+6\pi kc_1\int\operatorname{Tr}(dg_2g_2^{-1})(g_1^{-1}dg_1)
    \tag{4.28}
    \label{eq:4.28}
  \end{equation}
\end{enumerate}

\medskip\noindent\textbf{练习}\quad\emph{WZ 项的变分}

利用变分公式
\begin{equation}
  \delta(g^{-1}\partial_\mu g)=\partial_\mu(g^{-1}\delta g)+[g^{-1}\partial_\mu g,g^{-1}\delta g]
  \tag{4.29}
  \label{eq:4.29}
\end{equation}
我们计算：
\begin{equation}
  \frac{\partial}{\partial s}\operatorname{Tr}(g^{-1}dg)^{2n+1}
  =(2n+1)\operatorname{Tr}d\left(g^{-1}\frac{\partial g}{\partial s}\right)(g^{-1}dg)^{2n}
  \tag{4.30}
  \label{eq:4.30}
\end{equation}
含对易子的第二项消掉了。

现在与右端比较
\begin{equation}
  \begin{aligned}
    d\left[\operatorname{Tr}g^{-1}\frac{\partial g}{\partial s}(g^{-1}dg)^{2n}\right]
    ={}&\operatorname{Tr}d\left(g^{-1}\frac{\partial g}{\partial s}\right)(g^{-1}dg)^{2n}\\
    &+\operatorname{Tr}g^{-1}\frac{\partial g}{\partial s}\left[d\Theta\Theta^{2n-1}-\Theta d\Theta\Theta^{2n-2}\pm\cdots\right]
  \end{aligned}
  \tag{4.31}
  \label{eq:4.31}
\end{equation}
并利用 Maurer-Cartan 方程，我们发现第二组项成对抵消。$\spadesuit$
% ♠ 为原文的收尾记号（与开头的 \ensuremath{\clubsuit} 标记同一套作者手写符号），照原样保留。
